\documentclass[preview]{standalone}

\usepackage{amsmath}
\usepackage{amssymb}
\usepackage{stellar}
\usepackage{definitions}

\begin{document}

\id{il-cognitivismo}
\genpage

\section{Il cognitivismo}

\begin{snippetdefinition}{cognitivismo-definition}{Cognitivismo}
    Il \emph{cognitivismo} è un approccio psicologico che studia i processi
    attraverso cui il sistema cognitivo acquisisce, elabora, archivia e
    recupera le informazioni.
\end{snippetdefinition}

\begin{snippet}{nascita-cognitivismo}
    Il cognitivismo nasce dal superamento del comportamentismo. Un passaggio
    centrale è la critica di \emph{Noam Chomsky} a \emph{Verbal Behavior} di
    Skinner. Skinner interpreta l'acquisizione del linguaggio come risultato di
    processi ordinari di apprendimento; Chomsky sostiene invece che questa
    spiegazione riduce il linguaggio a ripetizioni stimolo-risposta e non rende
    conto della sua capacità combinatoria.
\end{snippet}

\begin{snippet}{competenza-innata-linguaggio}
    Secondo Chomsky, l'acquisizione del linguaggio implica una competenza
    innata: il bambino non si limita a ripetere espressioni apprese, ma può
    produrre e comprendere frasi nuove. Questa critica contribuisce alla
    nascita di una psicologia interessata ai processi interni di elaborazione
    dell'informazione.
\end{snippet}

\end{document}
